\chapter{Crossed products}\label{ch:crossed_products}

How adding a clock turns a type III algebra into a type II algebra: the crossed product, Takesaki duality, entropy and density matrices in type II, and finite-dimensional toy models.

\section{The crossed product construction}\label{ch:crossed}

\begin{physmotiv}
Modular flow is the hidden dynamics of a state, but on a type III algebra it is \emph{outer}: no operator of the algebra generates it (\cref{ch:modham}). The cure is to enlarge the algebra by the missing generator. Concretely, we adjoin a new degree of freedom (a ``clock'' $x$ with conjugate momentum $p$) and an operator $X=K+x$ whose adjoint action is the modular flow. In the enlarged algebra the flow is \emph{inner} by construction. This is the crossed product. It is a mathematical device in Part V; in gravity (Part VII) the clock is a physical observer and the enlargement is forced by the constraints.
\end{physmotiv}

\subsection{Abstract definition}

Let $\alpha:\R\to\Aut(\M)$ be a $\sigma$-weakly continuous one-parameter group of automorphisms of $\M\subset\BH$. On $L^2(\R,\HH)=\HH\otimes L^2(\R)$ define
\begin{equation}
(\pi_\alpha(a)\xi)(s)=\alpha_{-s}(a)\,\xi(s),\qquad(\lambda(t)\xi)(s)=\xi(s-t).
\end{equation}
Then $\pi_\alpha$ is a faithful representation of $\M$, $\lambda$ is a unitary representation of $\R$, and they satisfy the \emph{covariance relation}
\begin{equation}
\lambda(t)\,\pi_\alpha(a)\,\lambda(t)\ad=\pi_\alpha\big(\alpha_t(a)\big).
\end{equation}
(Check: both sides act on $\xi(s)$ by $\alpha_{t-s}(a)$.)

\begin{definition}
The \emph{crossed product}\index{crossed product} is the von Neumann algebra $\M\rtimes_\alpha\R=\{\pi_\alpha(a),\lambda(t):a\in\M,t\in\R\}''$ on $\HH\otimes L^2(\R)$.
\end{definition}

In words: $\M\rtimes_\alpha\R$ is generated by $\M$ and a unitary group $\lambda(t)$ that implements $\alpha_t$. It is the smallest algebra containing $\M$ in which $\alpha$ becomes inner. Elements have the form $\int\dd t\,a(t)\lambda(t)$ with $a(t)\in\M$, composed by the twisted convolution $(a\star b)(t)=\int\dd s\,a(s)\alpha_s(b(t-s))$.

\begin{figure}[h]
\centering
\begin{tikzpicture}[box/.style={draw,rounded corners,fill=physcol!7,align=center,font=\small,inner sep=6pt},node distance=3.2cm]
\node[box] (f1) {\textbf{Frame 1}\\ $\{a\otimes\id,\ X=K+x\}''$\\ $\Ad\ee^{-\ii tX}=\sigma_t$ on $\M$};
\node[box,right=4.2cm of f1] (f2) {\textbf{Frame 2}\\ $\{\sigma_p(a),\ x\}''$\\ commutes with $K-x$};
\draw[<->,thick] (f1)--(f2) node[midway,above,font=\small] {$\Ad\ee^{-\ii pK}$};
\node[box,below=1.6cm of f1] (fix) {$\M\otimes\id=N^{\theta}$\\ fixed points};
\draw[-{Stealth},thick] (f1)--(fix) node[midway,left,font=\small] {$\theta_s=\Ad\ee^{\ii sp}$};
\node[font=\small,align=center,right=0.3cm of fix] {$\theta_s:\ X\mapsto X+s,\quad a\mapsto a$};
\end{tikzpicture}
\sidecaption{The two descriptions of the crossed product $N=\M\rtimes_\sigma\R$. Frame~1 is convenient mathematically (the flow is inner, $\tau(f(X))=\int\ee^Xf(X)\,\dd X$); frame~2 is the physical one (dressed operators, invariants of the constraint). The dual action $\theta_s$ shifts $X$ and fixes $\M$.}
\end{figure}

\subsection{The physicist's form}

Take now $\alpha=\sigma^\Omega$, the modular flow of a cyclic separating vector $\Omega$. Put $K=-\log\Delta_\Omega$, so $\sigma_t(a)=\ee^{-\ii tK}a\,\ee^{\ii tK}$ and $K\Omega=0$. Let $x$ be the position operator on $L^2(\R)$ and $p=-\ii\,\dd/\dd x$.

\begin{proposition}
The crossed product is unitarily equivalent to
\begin{equation}
\M\rtimes_\sigma\R\;\cong\;\big\{\,a\otimes\id\ (a\in\M),\ \ X\equiv K+x\,\big\}''\ \subset\ \Bnd{\HH\otimes L^2(\R)}.
\end{equation}
Here $\ee^{-\ii tX}$ plays the role of $\lambda(t)$, and $X$ is affiliated with the algebra.
\end{proposition}
\emph{Check.} Since $x$ commutes with $a\otimes\id$, $\ee^{-\ii tX}(a\otimes\id)\ee^{\ii tX}=\ee^{-\ii tK}a\ee^{\ii tK}\otimes\id=\sigma_t(a)\otimes\id$. So $\ee^{-\ii tX}$ implements the modular flow, as required. (The identification with the $\pi_\alpha,\lambda$ picture is a unitary change of frame.)

\begin{keyidea}
In $\M\rtimes_\sigma\R$ the modular flow is \emph{inner}: it is $\Ad\,\ee^{-\ii tX}$ with $X=K+x$ in the algebra. The price is the new variable $x$. Note $K$ alone is \emph{not} in the algebra; only the combination $K+x$ is.
\end{keyidea}

\subsubsection*{Second frame (the one used in gravity)}
Conjugating by $\ee^{-\ii pK}$ and using $\ee^{-\ii pK}x\,\ee^{\ii pK}=x-K$ (since $\ee^{\ii cp}x\ee^{-\ii cp}=x+c$), we get $\ee^{-\ii pK}(K+x)\ee^{\ii pK}=x$ and
\begin{equation}
\M\rtimes_\sigma\R\;\cong\;\big\{\,\ee^{-\ii pK}a\,\ee^{\ii pK}\ (a\in\M),\ \ x\,\big\}''=\{\sigma_p(a),\,x\}'',\label{eq:frame2}
\end{equation}
where $\sigma_p(a)$ means the modular flow with the \emph{operator} $p$ as time. This is the form in which the algebra arises as the set of gauge-invariant observables in gravity (\cref{ch:add_observer,ch:clpw}): $x$ becomes (minus $\beta$ times) the observer's energy and $\sigma_p(a)$ is a field operator dressed to the observer. The two frames contain the same information; frame~1 is better for the mathematics, frame~2 for the physics.

\subsection{The commutant}

In frame 1 the commutant is again a crossed product:
\begin{equation}
(\M\rtimes_\sigma\R)'=\big\{\,\ee^{\ii pK}(a'\otimes\id)\ee^{-\ii pK}\ (a'\in\M'),\ \ x\,\big\}''.
\end{equation}
\emph{Check.} (i) $\ee^{\ii pK}a'\ee^{-\ii pK}=\Delta^{-\ii p}a'\Delta^{\ii p}\in\M'$ (the modular flow preserves $\M'$ too), so it commutes with $a\otimes\id$. (ii) $x$ commutes with $a\otimes\id$ and with $X=K+x$. (iii) Using $\ee^{-\ii pK}X\,\ee^{\ii pK}=x$, we get $[X,\ee^{\ii pK}a'\ee^{-\ii pK}]=\ee^{\ii pK}[x,a']\ee^{-\ii pK}=0$. That these generate the \emph{whole} commutant is part of Takesaki's theorem. In particular \emph{the commutant of the crossed product of the right wedge is (the crossed product of) the algebra of the left wedge}, again with the same variable $x$.

\subsection{Two sanity checks}

\subsubsection{Inner flows: the crossed product is trivial}
Let $\M=M_n$ on $\HH=\C^n\otimes\overline{\C^n}$ and $\Omega=\rho^{1/2}$, so $K=K_\M-K_{\M'}$ with $K_\M=-\log\rho\in\M$ (\cref{ch:tomita}). Then $y\equiv X-K_\M=x-K_{\M'}$ commutes with $\M$, with $K_\M$ and with $x$, and $X=K_\M+y$ with $K_\M\in\M$. Hence the algebra generated by $\M$ and $X$ is the algebra generated by $\M$ and $y$:
\begin{equation}
M_n\rtimes_{\sigma}\R\;\cong\;M_n\otimes\Linf(\R_y).
\end{equation}
Type I $\otimes$ abelian, with centre $\Linf(\R_y)$: the crossed product by an inner flow adds nothing but a classical variable. In particular it is of type I$_n$ over the centre and not a factor.

\subsubsection{The infinite spin chain}
For the Powers factor $R_\lambda$ with its periodic modular flow, $\sigma_{t_0}=\mathrm{id}$ for $t_0=2\pi/|\log\lambda|$. Then the unitary $\lambda(t_0)=\ee^{-\ii t_0X}$ commutes with $\M$ (because $\sigma_{t_0}=\mathrm{id}$) and with every $\lambda(t)$, so it lies in the centre; its spectrum is the circle dual to $t_0\Z$, and $\M\rtimes_\sigma\R$ is a direct integral over that circle of II$_\infty$ factors: type II$_\infty$ but \emph{not} a factor. For III$_1$ the flow has no period and the centre is trivial (next section).

\begin{whatwrong}
\begin{itemize}
\item \textbf{State independence.} The crossed product $\M\rtimes_{\sigma^\varphi}\R$ does not depend on which faithful normal state or weight $\varphi$ is used: for another $\psi$ the Connes cocycle $(D\psi:D\varphi)_t$ (\cref{ch:modham}) gives an isomorphism. Hence \emph{the algebra $\M\rtimes\R$ is canonical}, although its explicit realization uses a state.
\item \textbf{Not a tensor product.} $\M\rtimes_\sigma\R\not\cong\M\bar\otimes\Linf(\R)$ in general. The twisting by $\sigma$ is the whole point.
\item \textbf{Which time?} We cross by the \emph{full} modular flow. Crossing by a different one-parameter group (e.g.\ a flow that is not modular) gives different, possibly type III, algebras.
\end{itemize}
\end{whatwrong}

\subsection*{Exercises}
\begin{exercise}\label{ex:crossed:1}
Verify the covariance relation $\lambda(t)\pi_\alpha(a)\lambda(t)\ad=\pi_\alpha(\alpha_t(a))$ for $\pi_\alpha,\lambda$ as defined. Check that $\pi_\alpha$ is a $^*$-homomorphism.
\end{exercise}
\begin{exercise}\label{ex:crossed:2}
Verify $\ee^{\ii cp}x\,\ee^{-\ii cp}=x+c$ and the identity $\ee^{-\ii pK}(K+x)\ee^{\ii pK}=x$. Deduce \eqref{eq:frame2}.
\end{exercise}
\begin{exercise}\label{ex:crossed:3}
Verify that $x$ and $\ee^{\ii pK}a'\ee^{-\ii pK}$ commute with $a\otimes\id$ and with $K+x$.
\end{exercise}
\begin{exercise}\label{ex:crossed:4}
For $\M=M_2$ with $\rho=\mathrm{diag}(p,1-p)$, write out $K_\M$, $K_{\M'}$ and $y$ explicitly, and exhibit the isomorphism $\M\rtimes\R\cong M_2\otimes\Linf(\R_y)$. Which element is the centre?
\end{exercise}
\begin{exercise}\label{ex:crossed:5}
Let $\alpha_t=\Ad\,\ee^{-\ii tH}$ for a bounded self-adjoint $H\in\M$ (an inner flow), and let $X$ be the generator of $\lambda$ in frame 1 (so $\ee^{-\ii tX}$ implements $\alpha_t$). Show that $y=X-H$ commutes with $\M$ and with $X$, and conclude $\M\rtimes_\alpha\R\cong\M\bar\otimes\Linf(\R_y)$.
\end{exercise}

\section{Takesaki duality}\label{ch:takesaki}

\begin{physmotiv}
Having added a clock to a type III algebra, we want to know what we got. The answer is the structural theorem of the first half of the book: the crossed product of a type III$_1$ algebra by its modular flow is a type II$_\infty$ factor, it has a \emph{trace}\index{trace}, and the original algebra is recovered as the fixed points of a dual symmetry that rescales the trace. The existence of the trace is a consequence of the KMS property of the state, so it comes ``for free'' from thermal physics. The trace is the thing that will define entropy.
\end{physmotiv}

\subsection{The dual action}

On $N=\M\rtimes_\sigma\R$ (frame 1, \cref{ch:crossed}) define
\begin{equation}
\theta_s=\Ad\,\ee^{\ii sp},\qquad s\in\R.
\end{equation}
Since $\ee^{\ii sp}x\ee^{-\ii sp}=x+s$ and $p$ commutes with $a\otimes\id$ and with $K$, we have $\theta_s(a\otimes\id)=a\otimes\id$ and $\theta_s(X)=X+s$. So $\theta$ is a one-parameter group of automorphisms of $N$ that fixes $\M$ and translates $X$. In terms of the elements $\int\dd u\,a(u)\ee^{\ii uX}$, it multiplies the mode $u$ by $\ee^{\ii us}$.

\begin{proposition}
The fixed-point algebra is $N^\theta=\M\otimes\id\cong\M$.
\end{proposition}
(Intuitively: only the mode $u=0$ is invariant.) So we recover $\M$ from $(N,\theta)$.

\subsection{The dual weight and the trace}

Define the conditional expectation-like map $T:N_+\to\M_+$ (an \emph{operator-valued weight}\index{weight}) by averaging over the dual action, $T(b)=\int\dd s\,\theta_s(b)$ suitably normalized; on $b=\int\dd u\,a(u)\ee^{\ii uX}$ it picks out the zero mode $a(0)$. The \emph{dual weight} is
\begin{equation}
\hat\varphi=\omega\circ T,\qquad\omega=\inner\Omega{\,\cdot\,\Omega}.
\end{equation}
It is a faithful normal semifinite weight on $N$, $\theta$-invariant, and for a function of $X$ alone, $\hat\varphi(f(X))=\int\dd X\,f(X)$ (normalization fixed by convention). Its modular flow is $\sigma^{\hat\varphi}_t=\Ad\ee^{-\ii tX}$ (it restricts to $\sigma^\omega_t$ on $\M$ and fixes $X$).

\begin{theorem}[trace on the crossed product]\label{thm:trace}
Let $\tau(b)=\hat\varphi\big(\ee^{X/2}\,b\,\ee^{X/2}\big)$ for $b\in N_+$. Then $\tau$ is a faithful normal semifinite \emph{trace} on $N$, and it is rescaled by the dual action:
\begin{equation}
\tau\circ\theta_s=\ee^{-s}\,\tau.
\end{equation}
Moreover, for any bounded Borel function $f$,
\begin{equation}
\tau\big(f(X)\big)=\int_{-\infty}^\infty\dd X\;\ee^{X}f(X).
\end{equation}
\end{theorem}

\begin{proofsketch}
\emph{Trace property.} Write $b=\int\dd u\,a(u)\ee^{\ii uX}$. The weight $\hat\varphi$ has modular flow $\Ad\ee^{-\ii tX}$, so $\hat\varphi_h=\hat\varphi(h^{1/2}\,\cdot\,h^{1/2})$ has modular flow $\Ad(h^{\ii t})\circ\Ad\ee^{-\ii tX}$ (Connes cocycle, \cref{ch:modham}). This is trivial iff $h=\ee^{X}$, and a weight with trivial modular flow is a trace. Direct check of the same fact: with $b e^{X/2}=\int\dd u\,a(u)\ee^{(\ii u+1/2)X}$ one finds $\tau(b\ad b)=\int\dd u\,\omega\big(\tilde a(u)\ad\tilde a(u)\big)$ with $\tilde a(u)=a(u-\ii/2)$ and $\tau(bb\ad)=\int\dd u\,\omega\big(\tilde a(u)\,\sigma_{-\ii}(\tilde a(u)\ad)\big)$, and these agree exactly by the KMS condition $\omega(c\,\sigma_{-\ii}(d))=\omega(dc)$ (\cref{ch:tomita}). \emph{Scaling:} since $\hat\varphi$ is $\theta$-invariant and $\theta_s(\ee^{X/2})=\ee^{s/2}\ee^{X/2}$, $\tau\circ\theta_s=\ee^{-s}\tau$. The last formula follows from $\hat\varphi(f(X))=\int f\,\dd X$.
\end{proofsketch}

So the \emph{KMS property of $\omega$ for its own modular flow is what makes $\tau$ tracial}. This is the mathematical heart of the whole book.

\begin{whatwrong}
Conventions matter and are easily misread across papers: (i) which sign of $\ii$ in $\Delta^{\pm\ii t}$; (ii) whether the dual action is $\Ad\ee^{\ii sp}$ or its inverse, i.e.\ whether the trace scales as $\ee^{-s}$ or $\ee^{+s}$; (iii) whether the variable called $x$ is $X$ or $-X$. We have fixed them in \cref{ch:crossed} (see also the style guide) and checked them in finite dimensions below; the physics statements (existence of the trace, finiteness for $x\le0$) are convention independent.
\end{whatwrong}

\subsection{Type of $N$ and Takesaki duality}

\begin{theorem}[Connes--Takesaki, as stated in \cref{ch:connes_class}]
If $\M$ is a $\sigma$-finite type III$_1$ factor, then $N=\M\rtimes_\sigma\R$ is a type II$_\infty$ factor with trace $\tau$ above. In general (type III, $\sigma$-finite), $N$ is semifinite: it has a f.n.s.\ trace.
\end{theorem}
The factor property uses that $\sigma$ has no ``periodic part'' (the centre of $N$ is the \emph{flow of weights}, which is trivial for III$_1$). Also:
\begin{theorem}[Takesaki duality]
$N\rtimes_\theta\R\;\cong\;\M\bar\otimes\Bnd{L^2(\R)}$.
\end{theorem}
So crossing twice returns the original algebra, up to tensoring with type I$_\infty$: the type is unchanged. Since $\M\otimes\Bnd{L^2}$ and $\M$ are Morita equivalent, one can say $\M$ is recovered from $(N,\tau,\theta)$: \emph{a type III$_1$ factor is the same data as a II$_\infty$ factor with a trace-scaling flow}.

\subsubsection*{II$_\infty$ versus II$_1$}
Because $\tau\circ\theta_s=\ee^{-s}\tau$, we have $\tau(\id)=\infty$ (indeed $\int\ee^{X}\dd X=\infty$). A type II$_\infty$ algebra has finite-trace \emph{projections}, and one can restrict to a corner. For $P=\chi(X\le0)\in N$,
\begin{equation}
\tau(P)=\int_{-\infty}^0\ee^{X}\dd X=1.
\end{equation}
The corner $PNP$ is a type II$_1$ factor with normalized trace $\tau(P\,\cdot\,P)$. In gravity the condition $X\le0$ will be the statement that the observer's energy is non-negative (\cref{ch:add_observer}); the maximum-entropy state will then have density matrix $P$ (\cref{ch:entropy_II,ch:clpw}).

\subsection{Finite-dimensional check}
For $\M=M_n$ (\cref{ch:crossed}), $N=M_n\otimes\Linf(\R_y)$ with centre variable $y=X+\log\rho$. A direct computation gives
\begin{equation}
\tau(F)=\int\dd y\,\ee^{y}\,\Tr F(y),\qquad F(y)\in M_n,
\end{equation}
and for $F=f(X)$, whose eigenvalues are $f(y-\log p_i)$ (since $X=y-\log\rho$),
\begin{equation}
\tau\big(f(X)\big)=\sum_i\int\dd y\,\ee^{y}f(y-\log p_i)=\sum_ip_i\int\dd X\,\ee^{X}f(X)=\int\dd X\,\ee^Xf(X),
\end{equation}
using $y=X+\log p_i$ and $\sum_ip_i=1$. This agrees with \cref{thm:trace}, and $\tau\circ\theta_s=\ee^{-s}\tau$ follows from $\theta_s:y\mapsto y+s$. Of course type I algebras already have traces; the point is that the formula $\tau(f(X))=\int\ee^Xf(X)\dd X$ is the one that survives when $\rho$ ceases to exist.

\begin{keyidea}
Cross the type III$_1$ factor by its modular flow: the result is a type II$_\infty$ factor with trace $\tau(f(X))=\int\ee^{X}f(X)\dd X$ coming from the KMS property. A dual flow $\theta_s$ rescales the trace by $\ee^{-s}$ and its fixed points are the original algebra. Cutting down to $X\le0$ gives a type II$_1$ corner.
\end{keyidea}

\subsection*{Exercises}
\begin{exercise}\label{ex:takesaki:1}
Show that the fixed points of $\theta_s=\Ad\ee^{\ii sp}$ on $\{a\otimes\id,K+x\}''$ are $\M\otimes\id$: argue with the Fourier decomposition $b=\int\dd u\,a(u)\ee^{\ii uX}$.
\end{exercise}
\begin{exercise}\label{ex:takesaki:2}
Verify the finite-dimensional trace formula. Use $N=M_n\otimes\Linf(\R_y)$ with $y=X+\log\rho$ and check (a) $\tau(FG)=\tau(GF)$ for $F,G\in N$ (trivial: $\Tr$ cyclic), (b) $\tau\circ\theta_s=\ee^{-s}\tau$ with $\theta_s:y\mapsto y+s$, (c) $\tau(f(X))=\int\ee^Xf(X)\dd X$ using $\sum p_i=1$.
\end{exercise}
\begin{exercise}\label{ex:takesaki:3}
Show that $\tau(\chi(X\le0))=1$ and $\tau(\chi(X\ge0))=\infty$. Compute $\tau(\chi(X\le c))$ and show that all values in $(0,\infty)$ occur, so $N$ has projections of every dimension (type II).
\end{exercise}
\begin{exercise}\label{ex:takesaki:4}
Verify the claim in the proof sketch: a weight $\varphi_h=\varphi(h^{1/2}\cdot h^{1/2})$ has $\sigma^{\varphi_h}_t=\Ad(h^{it})\circ\sigma^\varphi_t$ in finite dimensions ($\varphi=\Tr(\rho\,\cdot)$, $\varphi_h=\Tr(h^{1/2}\rho h^{1/2}\,\cdot)$ with $h$ commuting). Why is this trivial iff $h\propto\rho^{-1}$?
\end{exercise}

\section{Entropy in type II}\label{ch:entropy_II}

\begin{physmotiv}
The crossed product is type II$_\infty$, with a trace $\tau$ (\cref{ch:takesaki}). So it has density matrices $\hat\rho$ and a von Neumann entropy $S=-\tau(\hat\rho\log\hat\rho)$. What does this entropy look like for states built from the original QFT state plus a clock? This section answers the question \emph{exactly} in the finite-dimensional model, where everything can be computed, and then states what survives when the underlying algebra is type III. The answer has the shape of a generalized-entropy formula: minus a relative entropy plus the entropy and energy of the clock.
\end{physmotiv}

\subsection{Density matrices and entropy in a semifinite algebra}

Let $N$ have a f.n.s.\ trace $\tau$ (\cref{ch:traces}). A normal state $\hat\Psi$ has a unique density matrix $\hat\rho\ge0$ affiliated with $N$, $\tau(\hat\rho)=1$, $\hat\Psi(b)=\tau(\hat\rho\,b)$, and
\begin{equation}
S(\hat\Psi)=-\tau(\hat\rho\log\hat\rho).
\end{equation}
Rescaling $\tau\to c\,\tau$ shifts $S$ by $\log c$ for every state: \emph{only differences of entropies are meaningful until a normalization is fixed.} In II$_1$ the natural normalization is $\tau(\id)=1$, which sets the maximum to $0$. In II$_\infty$ with no finite $\tau(\id)$ the zero of entropy is conventional, but by \cref{ch:takesaki} the dual action fixes the scaling $\tau\circ\theta_s=\ee^{-s}\tau$, and the corner $P=\chi(X\le0)$ has $\tau(P)=1$.

\subsection{States of the system plus clock}

Work in frame 1 of \cref{ch:crossed}. Let $\Phi\in\HH$ be a vector (a state of the original system) and $g\ge0$ a probability density on $\R$. The \emph{clock-dressed state}\index{clock} is the vector
\begin{equation}
\hat\Phi=\Phi\otimes g^{1/2}(x)\in\HH\otimes L^2(\R),
\end{equation}
regarded as a state of $N$. The function $g$ is the probability distribution of the clock variable $x$. We compute the entropy of $\hat\Phi$ exactly in the finite-dimensional model.

\subsubsection{Exact finite-dimensional computation}
Let $\M=M_n$ with $\Omega=\rho^{1/2}=\sum_i\sqrt{p_i}\ket i\ket{\bar i}$, so $K=K_\M-K_{\M'}$ and $N=M_n\otimes\Linf(\R_y)$ with centre variable $y=X-K_\M=x-K_{\M'}$ (\cref{ch:crossed}). On the basis vector $\ket i\ket{\bar i}\otimes\ket x$ the operator $y$ acts as multiplication by $x+\log p_i$. Take $\Phi=u\Omega$ with $u\in\M$ a unitary (a ``locally excited'' state; $u=\id$ gives $\Omega$). Then $\Phi=\sum_i\sqrt{p_i}\ket{ui}\ket{\bar i}$ and an element $F=\sum_if_i(y)\ket{ui}\bra{ui}$ of $N$ has
\begin{equation}
\begin{split}\hat\Phi(F)&=\sum_ip_i\int\dd x\,g(x)f_i(x+\log p_i)=\int\dd y\,\ee^{y}\sum_i\hat\rho_i(y)f_i(y),\\ \hat\rho_i(y)&=\ee^{-y}p_i\,g(y-\log p_i),\end{split}
\end{equation}
so the density matrix has eigenvalues $\hat\rho_i(y)$ on the eigenvectors $\ket{ui}$ (check $\tau(\hat\rho)=\sum_ip_i\int g=1$). Therefore
\begin{align}
S(\hat\Phi)&=-\sum_i\int\dd y\,\ee^y\hat\rho_i\log\hat\rho_i
=-\sum_ip_i\int\dd x\,g(x)\big[-x+\log g(x)\big]\notag\\
&=\langle x\rangle_g+h(g),\qquad h(g)=-\int g\log g\,\dd x.\label{eq:Sexact}
\end{align}
(We used $y=x+\log p_i$, so $-y+\log p_i=-x$; the $p_i$ have dropped out.) The entropy depends only on the clock distribution. We checked this formula numerically for several $p_i$ and Gaussian $g$ (\texttt{scripts/check\_toy.py}). To express it intrinsically, note $\langle X\rangle_{\hat\Phi}=\langle K\rangle_\Phi+\langle x\rangle_g$, and $\langle K\rangle_{u\Omega}=S_{\rm rel}(u\Omega\|\Omega)$ by the exact identity of \cref{ch:modham}. Hence
\begin{equation}
\boxed{\ S(\hat\Phi)=h(g)+\langle X\rangle_{\hat\Phi}-S_{\rm rel}(\Phi\|\Omega)\ }\label{eq:Srel}
\end{equation}
Equivalently: entropy $=$ clock entropy $+$ clock-plus-modular ``energy'' $-$ relative entropy of the system state.

\begin{keyidea}
In the crossed product the entropy of a clock-dressed state is, up to the clock terms, \emph{minus the relative entropy}\index{relative entropy} of the system's state with respect to the reference state $\Omega$. The divergent entropy of QFT has been replaced by $-S_{\rm rel}$, which is finite, plus the clock contribution $h(g)+\langle X\rangle$.
\end{keyidea}

\subsection{The maximum entropy state}

Take $\Phi=\Omega$ and $g_{\max}(x)=\ee^{x}\,\theta(-x)$, a normalized density ($\int_{-\infty}^0\ee^x\dd x=1$). Then $\hat\rho_i(y)=\ee^{-y}p_i\ee^{y-\log p_i}\theta(-(y-\log p_i))=\theta(\log p_i-y)$, which is exactly the spectral projection $P=\chi(X\le0)$ (since $X$ has eigenvalue $y-\log p_i$ on the $i$th sector). So
\begin{equation}
\hat\rho_{\max}=P=\chi(X\le0),\qquad S=-\tau(P\log P)=0.
\end{equation}
By the Jensen argument of \cref{ch:traces}, every normal state of the II$_1$ corner $PNP$ has $S\le0=S(\hat\rho_{\max})$: \emph{this is the maximum-entropy state}\index{maximum-entropy state}, the ``infinite-temperature'' state of $PNP$ (its density matrix is the identity of the corner). The same value follows from \eqref{eq:Sexact}: $\langle x\rangle_{g_{\max}}=-1$, $h(g_{\max})=-\int\ee^xx\dd x=+1$, sum zero.

More generally, for $\Phi=u\Omega$ and any $g$ supported on $x\le0$:
\begin{equation}
S(\hat\Phi)=-\int g\log\frac{g}{g_{\max}}=-D_{\rm KL}(g\,\|\,g_{\max})\le0.
\end{equation}
The clock's distribution relative to $\ee^{x}$ is a Kullback--Leibler divergence: the maximum is attained when the clock is distributed as $\ee^{x}$. Looking ahead (\cref{ch:add_observer}): $x=-\beta q$ with $q\ge0$ the observer's energy, so $g_{\max}\propto\ee^{-\beta q}$ is a \emph{Boltzmann distribution} at the horizon temperature: the maximum entropy state has the observer in thermal equilibrium with the horizon.

\subsection{What survives for a type III algebra}

For a type III$_1$ algebra $\M$ no density matrix $\rho$ exists, so the derivation above (which used $\rho=\sum p_i\ket i\bra i$) does not apply literally; but its \emph{result} is formulated using only objects that exist: $\Omega$, $K$, $S_{\rm rel}$, $X$, $g$. The theorem is that in the semiclassical regime, where $g$ is a slowly varying function of $x$ and $\Phi$ is a state with $\langle K\rangle_\Phi$ finite and small compared to the scale of variation of $g$, the density matrix of $\Phi\otimes g^{1/2}$ in the type II$_\infty$ algebra is, up to corrections that vanish in this limit,
\begin{equation}
S(\hat\Phi)=h(g)+\langle X\rangle_{\hat\Phi}-S_{\rm rel}(\Phi\|\Omega)+(\text{corrections vanishing in this limit}),
\end{equation}
that is, \eqref{eq:Srel} survives, with $S_{\rm rel}$ now Araki's relative entropy of \cref{ch:modham}. (The explicit density matrix in the type III setting is constructed in the original papers; we only need its entropy.) We shall see in \cref{ch:anatomy} how this reproduces the generalized entropy $\Sgen=A/4G+S_{\rm bulk}$ when the clock is an observer in de Sitter space. The precise statement (a theorem of Chandrasekaran--Longo--Penington--Witten in this setting) requires the regime of validity to be spelled out; we do so there and flag the source-checking in the bibliography.

\begin{whatwrong}
\begin{itemize}
\item The entropy \eqref{eq:Sexact} is \emph{not} the entropy of the original system. It is the entropy of the algebra $N$ in the state $\hat\Phi$; its dependence on $\Phi$ is only through the relative entropy to $\Omega$ and through $\langle K\rangle$.
\item The cutoff dependence of the QFT entropy has not vanished by magic: it has moved into the \emph{choice of trace normalization}, i.e.\ a state-independent additive constant. Differences of entropies between states are unambiguous.
\item In the finite-dimensional model all states of $N$ are covered by general $\hat\rho$; ours ($\Phi\otimes g^{1/2}$) is a special family, and \eqref{eq:Srel} was derived only for $\Phi=u\Omega$.
\end{itemize}
\end{whatwrong}

\subsection*{Exercises}
\begin{exercise}\label{ex:entropy_II:1}
Verify \eqref{eq:Sexact}: from the formula for $\hat\rho_i(y)$ compute $-\int\dd y\,\ee^y\sum_i\hat\rho_i\log\hat\rho_i$ and show the $p_i$-dependence drops out. Run \texttt{scripts/check\_toy.py} and compare.
\end{exercise}
\begin{exercise}\label{ex:entropy_II:2}
For $g$ a Gaussian with mean $\mu$ and variance $s^2$ compute $h(g)+\langle x\rangle$. For which $(\mu,s)$ is $S=0$? Show $S\le0$ cannot hold for all Gaussians, and explain why: Gaussians have support on $x>0$ where $g_{\max}=0$.
\end{exercise}
\begin{exercise}\label{ex:entropy_II:3}
Prove $S=-D_{\rm KL}(g\|g_{\max})$ for $g$ supported on $x\le0$ and conclude $S\le0$ with equality iff $g=g_{\max}$. Check with $g(x)=\lambda\ee^{\lambda x}\theta(-x)$: show $S=1-\frac1\lambda-\log\lambda$ and that its maximum over $\lambda>0$ is $0$ at $\lambda=1$.
\end{exercise}
\begin{exercise}\label{ex:entropy_II:4}
Show that rescaling $\tau\to c\tau$ changes the density matrix $\hat\rho\to\hat\rho/c$ and $S\to S+\log c$, and relate this to the choice $\tau(P)=1$.
\end{exercise}

\section{Toy models: a system plus a clock}\label{ch:toy_clock}

\begin{physmotiv}
Before touching fields and gravity, here is the entire construction in a model small enough to hold in your head: a qubit with a thermal state, plus a clock. The model is type I, so nothing is truly type III, but \emph{all} the structural features of the de Sitter construction appear: the constraint, the dressed operators, the trace with weight $\ee^{x}$, the half-line $x\le0$ that makes the trace finite, and the Boltzmann-distributed clock in the maximum-entropy state. If this section makes sense, \cref{ch:clpw} is the same story with $\HH$ replaced by Fock space.
\end{physmotiv}

\subsection{The model}

A system (``static patch'') is a qubit: $\M=M_2$ with Hamiltonian $H=\mathrm{diag}(E_0,E_1)$, level splitting $\omega_0=E_1-E_0>0$, and a thermal state $\rho=\ee^{-\beta H}/Z=\mathrm{diag}(p_0,p_1)$. Purify it by the thermofield double
\begin{equation}
\Omega=\sqrt{p_0}\ket{0}\ket{\bar0}+\sqrt{p_1}\ket1\ket{\bar1}\in\C^2\otimes\overline{\C^2},
\end{equation}
so that $\M=M_2\otimes\id$ is the algebra of the ``right patch'' and $\M'$ the algebra of the ``left''. The modular Hamiltonian is $K=\beta(H_R-H_L)$ (\cref{ch:tomita}) and the modular flow acts on matrix units by
\begin{equation}
\sigma_t\big(\ket i\!\bra j\big)=\Big(\frac{p_i}{p_j}\Big)^{\ii t}\ket i\!\bra j=\ee^{-\ii t\beta(E_i-E_j)}\ket i\!\bra j.
\end{equation}

\subsection{The constraint and the dressed operators}

Add a clock: a particle on a line with position $x$ and momentum $p$, so the total Hilbert space is $\HH\otimes L^2(\R)$. Impose a gauge constraint that ties the clock to the system's modular time: physical (gauge-invariant) operators are those invariant under
\begin{equation}
\alpha_t=\sigma_t\otimes\Ad\,\ee^{\ii tx}=\Ad\,\ee^{-\ii t(K-x)},
\end{equation}
that is, they commute with the constraint generator $\hat G=K-x$ (\cref{ch:crossed}, frame 2; we identify $K=\beta H$ and $x=-\beta q$ so that $\hat G=\beta(H+q)$, the total energy of the system plus the clock-observer, see \cref{ch:add_observer}). Using $\ee^{\ii tx}p\,\ee^{-\ii tx}=p-t$ one checks that the \emph{dressed operators}
\begin{equation}
\hat a=\sigma_p(a)=\ee^{-\ii pK}a\,\ee^{\ii pK},\qquad x
\end{equation}
are invariant, and they generate $N=\{\sigma_p(a),x\}''\cong\M\rtimes_\sigma\R$.

On matrix units: $\widehat{\ket i\!\bra j}=\ket i\!\bra j\,\ee^{-\ii\beta(E_i-E_j)p}$. Since $\ee^{-\ii cp}$ maps $\ket x\mapsto\ket{x+c}$, the dressed operator $\widehat{\ket1\!\bra0}$ raises the system energy by $\omega_0$ and simultaneously shifts $x$ by $c=\beta\omega_0$, i.e.\ shifts the clock energy $q=-x/\beta$ by $-\omega_0$. \emph{Energy conservation, $H+q=$ const, is enforced operator by operator:} a system transition is physical only if accompanied by a compensating change in the clock.

\begin{keyidea}
Gauge-invariant operators of the system are the system operators dressed by $\ee^{-\ii p\beta(E_i-E_j)}$: each transition is compensated by a shift of the clock. The algebra of invariants is the crossed product.
\end{keyidea}

\subsection{The algebra, trace and entropy}

By \cref{ch:crossed}, $N\cong M_2\otimes\Linf(\R_y)$, type I over its centre $y=X-K_\M$. The trace is $\tau(F)=\int\dd y\,\ee^{y}\Tr F(y)$ (\cref{ch:takesaki}). \emph{Every} state of the form $\hat\Phi=u\Omega\otimes g^{1/2}(x)$ has entropy
\begin{equation}
S(\hat\Phi)=\langle x\rangle_g+h(g)\qquad(\text{\cref{ch:entropy_II}}),
\end{equation}
independent of $\beta$, $E_i$ and $u$. With $x=-\beta q$ and $g(x)\dd x=f(q)\dd q$ this reads
\begin{equation}
S(\hat\Phi)=-\beta\langle q\rangle_f+h(f)+\log\beta.
\end{equation}
(Here $g(x)\dd x=f(q)\dd q$ gives $h(g)=h(f)+\log\beta$; the $\log\beta$ is the kind of additive constant fixed by the trace normalization.) Up to this constant, this is the thermodynamic functional $-\beta\langle q\rangle+h(f)$, maximized subject to $q\ge0$ by $f(q)=\beta\ee^{-\beta q}$, the Boltzmann distribution (indeed $\langle q\rangle=1/\beta$, $h(f)=1-\log\beta$, so $S=-1+1-\log\beta+\log\beta=0$). At the maximum $S=0$ with $\tau(P)=1$ (\cref{ch:entropy_II}).

\subsubsection*{Why the half-line $x\le0$ matters}
\begin{center}
\renewcommand{\arraystretch}{1.3}
\resizebox{\ifdim\width>\linewidth\linewidth\else\width\fi}{!}{%
\begin{tabular}{@{}lll@{}}
\toprule
Clock spectrum & $\tau(\id)$ & Algebra\\
\midrule
$q\in\R$ (all $x\in\R$) & $\int_{-\infty}^\infty\ee^{x}\dd x=\infty$ & II$_\infty$; no maximum entropy state\\
$q\ge0$ ($x\le0$) & $\int_{-\infty}^0\ee^{x}\dd x=1$ & II$_1$; maximum entropy state $P$\\
\bottomrule
\end{tabular}}
\end{center}
This is the same elementary fact as in \cref{ch:groups}: $\int_0^\infty\ee^{-\beta q}\dd q$ converges, $\int_{-\infty}^\infty\ee^{-\beta q}\dd q$ does not. A clock with Hamiltonian bounded below is what turns II$_\infty$ into II$_1$.

\subsection{What the toy teaches and what it does not}

\textbf{Teaches.}
(i) The constraint $H+q$ produces the crossed product as its algebra of invariants. (ii) The trace has weight $\ee^{x}$ and its finiteness depends on the clock spectrum. (iii) The maximum-entropy state is the original state (here: thermofield double) with the clock in a Boltzmann distribution. (iv) Entropies of dressed states are $-S_{\rm rel}$ plus clock terms.

\textbf{Does not.}
(a) Nothing here is type III: the divergences are absent because the system has finitely many levels, and the entropy $\langle x\rangle+h(g)$ is independent of the system state, because for a type I algebra the dependence on $\rho$ cancels. In the genuine QFT/gravity case the \emph{regularized} bulk entropy plus the area term combine into $\Sgen$, and the relative entropy dependence is visible (\cref{ch:anatomy}). (b) In the toy the observer is a decoupled degree of freedom, with no back-reaction; in gravity the observer's mass sets the horizon size and the constraint involves the full dS isometry group (\cref{ch:ds_special,ch:what_observer}). (c) The toy has no notion of a semiclassical limit; the real construction needs $\beta q\gg1$ states.

\subsection*{Exercises}
\begin{exercise}\label{ex:toy_clock:1}
Compute $\widehat{\ket1\!\bra0}$ explicitly and verify that it commutes with $\hat G=K-x$ (on $\HH\otimes L^2$). Check that $\ket1\!\bra0\otimes\id$ alone does not.
\end{exercise}
\begin{exercise}\label{ex:toy_clock:2}
For the qubit, list the generators of $N$ and exhibit $N\cong M_2\otimes\Linf(\R_y)$ concretely: find the central element $y$ in terms of $x$ and $K_{\M'}$, and the matrix units $e_{ij}\in\M$.
\end{exercise}
\begin{exercise}\label{ex:toy_clock:3}
Using \cref{ch:entropy_II}, maximize $S=-\beta\langle q\rangle+h(f)+\log\beta$ over densities $f$ on $q\ge0$ (here $h(f)=-\int f\log f\,\dd q$ is taken in the variable $q$; the $\log\beta$ comes from $x=-\beta q$, $g=f/\beta$) and show $f=\beta\ee^{-\beta q}$ with maximum $S=0$. Compute $\langle q\rangle$ and interpret it as the mean energy of a classical thermal oscillator mode at the temperature $1/\beta$.
\end{exercise}
\begin{exercise}\label{ex:toy_clock:4}
Extend to a three-level system with energies $E_0,E_1,E_2$. List all gauge-invariant dressed operators and the energy shifts of the clock they induce. Show the entropy of $\Omega\otimes g^{1/2}$ is unchanged.
\end{exercise}
\begin{exercise}\label{ex:toy_clock:5}
(Simulation.) Discretize $x$ on a grid and verify numerically $\tau(\hat\rho)=1$ and $S=0$ for the maximum entropy state, modifying \texttt{scripts/check\_toy.py}. Show $S<0$ for any other density supported on $x\le0$.
\end{exercise}

